%-------------------------------------------------------------------------------
%	SECTION TITLE
%-------------------------------------------------------------------------------
\cvsection{Personal Value Statement}


%-------------------------------------------------------------------------------
%	CONTENT
%-------------------------------------------------------------------------------
\begin{cvparagraph}

%---------------------------------------------------------
我的核心竞争力在于\textbf{\textcolor{awesome}{生物信息算法的原创开发能力}}：不仅能熟练运用现有分析流程，更能在\textbf{\textcolor{awesome}{现有工具无法满足研究需求时，独立设计并实现原创算法}}。这一能力让我可以为不同研究方向\textbf{\textcolor{awesome}{提供定制化的生物信息解决方案}}，成为多个课题组的共享生信支撑力量。

在博士阶段，我独立开发了两款代表性的生物信息工具：\textbf{\textcolor{awesome}{BALEEN}}——基于\textbf{动态时间规整}与\textbf{贝叶斯模型}的Nanopore直接RNA测序修饰检测算法（Nature Methods在投）；\textbf{\textcolor{awesome}{TAPE}}——基于\textbf{深度自编码器}的单细胞反卷积工具，已发表于\textbf{Nature Communications}（IF: 15.7）。这些方法已为多个香港RGC GRF获批项目及联合基金项目提供关键技术支撑，充分验证了\textbf{\textcolor{awesome}{原创算法解决真实科研问题}}的能力。

我具备\textbf{\textcolor{awesome}{全栈生信开发能力}}：底层熟悉C/CUDA高性能计算与信号处理；中层精通Python/R的数据分析与可视化；模型层掌握PyTorch深度学习框架（自编码器、Transformer等架构）与传统机器学习/统计建模；工程层熟练运用Docker、Singularity、Git等工具构建可复现、可部署的分析流程。这一技术栈可\textbf{\textcolor{awesome}{支撑从单分子测序到多组学整合}}的各类分析任务，涵盖NGS基因组/转录组、单细胞转录组、Nanopore长读长测序、宏基因组、比较基因组学、种群遗传学等主流方向。

通过搭建\textbf{\textcolor{awesome}{生物信息学公共服务平台}}，我可以系统性地提升各课题组的科研产出效率——从实验设计、数据分析到论文撰写提供\textbf{全流程支持}，整合分散的研究力量、凝练共同研究方向，为学科建设与硕士点申报提供成果支撑。同时，我具备\textbf{\textcolor{awesome}{基金申请的实战经验}}（主要参与3项获批基金，总计405万港币），可协助团队积极申报国家自然科学基金及省部级课题。

此外，我拥有\textbf{4年业界研发与管理经验}（含生物信息总监），主导过多个\textbf{\textcolor{awesome}{临床检测平台与全自动分析系统的从0到1建设}}，授权发明专利5项，具备团队组建、项目统筹与技术落地能力，能够快速建立高效的科研团队与研究生培养体系。

\end{cvparagraph}
